% !TEX root = IE3_expC_prelab.tex
\documentclass[lualatex,a4paper,10pt]{jlreq}
% Shared LuaLaTeX preamble.
\usepackage{luatexja-fontspec}
\setmainfont{TeX Gyre Termes}
\setsansfont{TeX Gyre Heros}
\setmonofont{Latin Modern Mono}
\setmainjfont[BoldFont=HaranoAjiGothic-Medium]{HaranoAjiMincho-Regular}
\setsansjfont[BoldFont=HaranoAjiGothic-Bold]{HaranoAjiGothic-Medium}
\usepackage[top=22mm,bottom=22mm,left=23mm,right=23mm]{geometry}
\usepackage{amsmath,amssymb,booktabs,array,longtable,tabularx}
\usepackage{graphicx,xcolor,tikz,circuitikz}
\usetikzlibrary{arrows.meta,positioning,calc}
\usepackage{enumitem,fancyhdr,listings,needspace}
\usepackage[unicode,colorlinks=true,linkcolor=labblue,urlcolor=labteal]{hyperref}
\usepackage{bookmark}
\definecolor{labblue}{RGB}{30,70,112}
\definecolor{labteal}{RGB}{0,100,105}
\definecolor{labgray}{gray}{0.95}
\ModifyHeading{section}{font={\large\sffamily\bfseries\color{labblue}}}
\ModifyHeading{subsection}{font={\normalsize\sffamily\bfseries\color{labteal}}}
\setlength{\parindent}{1\zw}
\setlength{\parskip}{3pt}
\setlength{\emergencystretch}{2em}
\renewcommand{\arraystretch}{1.35}
\setlist{itemsep=3pt,topsep=4pt,leftmargin=2\zw}
\newcolumntype{Y}{>{\raggedright\arraybackslash}X}
\newcolumntype{C}[1]{>{\centering\arraybackslash}p{#1}}
\pagestyle{fancy}
\fancyhf{}
\fancyhead[L]{\small\color{labblue} IE3 後期・電子工学実験}
\fancyhead[R]{\small テーマ\LabCode}
\fancyfoot[C]{\thepage}
\setlength{\headheight}{16pt}
\DeclareOldFontCommand{\rm}{\normalfont\rmfamily}{\mathrm}
\lstset{basicstyle=\small\ttfamily,columns=fullflexible,keepspaces=true,
 breaklines=true,frame=single,backgroundcolor=\color{labgray},
 numbers=left,numberstyle=\tiny,showstringspaces=false,aboveskip=8pt,belowskip=8pt}
\newcommand{\blank}{\rule{22mm}{0.3pt}}
\newcommand{\answerline}{\par\noindent\rule{\linewidth}{0.25pt}\par}
\newcommand{\writing}[1]{\par\Needspace{\dimexpr#1+5mm\relax}\noindent%
 \fcolorbox{labblue!35}{white}{\parbox[c][#1][t]{\dimexpr\linewidth-2\fboxsep-2\fboxrule\relax}{\vspace{2mm}\noindent\strut\hfill\par}}\par}
\newcommand{\hint}[1]{\par\smallskip\noindent{\sffamily\bfseries\color{labteal}記述の視点：}#1\par}
\newcommand{\note}[1]{\par\smallskip\noindent\fcolorbox{labblue!45}{labblue!5}{%
 \parbox{\dimexpr\linewidth-2\fboxsep-2\fboxrule\relax}{#1}}\par\smallskip}
\newcommand{\eqerror}{\varepsilon=\frac{x_{\mathrm{meas}}-x_{\mathrm{th}}}{x_{\mathrm{th}}}\times100\ \%}
\newcommand{\LabSource}{徳山工業高等専門学校 情報電子工学科，
 『2026年度 電子工学実験（3年後期）テキスト』，テーマ\LabCode，
 \expandafter\path\expandafter{\SourcePdf}。}
\newcommand{\reporttitle}{
 \noindent{\small\sffamily\color{labteal}実験報告書・記入ガイド}\par
 \noindent{\LARGE\sffamily\bfseries \LabCode\quad\LabTitle}\par\smallskip
 \noindent 氏名：\blank\quad 班：\blank\quad 実験日：\blank\par
 \note{目的・原理・手順を確認し、実測値と自分の考察を記入する。
 考察のガイドは、検討する事象・使う測定結果・記述の方針を示す。}
}
\newcommand{\prelabtitle}{
 \noindent{\small\sffamily\color{labteal}事前学習・前提知識プリント}\par
 \noindent{\LARGE\sffamily\bfseries \LabCode\quad\LabTitle}\par\smallskip
 \noindent 氏名：\blank\quad 班：\blank\quad 予習日：\blank\par
 \note{用語、回路の動き、計算、測定表への記入の順に読む。
 例題の値は説明用であり、実験結果ではない。}
}
\newcommand{\tocrow}[3]{\hyperref[#1]{#2}& #3 &\pageref*{#1}\\[2mm]}
\newcommand{\documentmap}[1]{
 \section*{目次と概要}
 \noindent 青色の項目をクリックすると該当箇所へ移動する。\par
 {\small\begin{longtable}{@{}p{0.29\linewidth}p{0.61\linewidth}r@{}}
 \toprule {\color{labblue}\bfseries 項目} & {\color{labblue}\bfseries 読むと分かること} & 頁\\\midrule
 \endhead #1\bottomrule
 \end{longtable}}
 \clearpage
}
\newenvironment{terms}{\par\smallskip\begin{longtable}{@{}p{0.23\linewidth}p{0.73\linewidth}@{}}
 \toprule {\color{labteal}\bfseries 用語}&{\color{labteal}\bfseries 意味・回路での役割}\\\midrule\endhead}
 {\bottomrule\end{longtable}}
\newcommand{\term}[2]{\textbf{#1}&#2\\[2mm]}
\newcommand{\guidepoint}[4]{
 \Needspace{32mm}\subsection{#1}
 \noindent{\bfseries\color{labteal}考える事象：}#2\par
 \noindent{\bfseries\color{labteal}参照する結果：}#3\par
 \noindent{\bfseries\color{labteal}記述の方針：}#4\par
}
\newcommand{\reportclosing}[1]{
 \clearpage\section{結論のまとめ方}\label{sec:closing}
 \hint{#1}
 \ConclusionGuide
 \begin{enumerate}
 \item 目的に対応する最も重要な実測結果を、測定条件と数値を添えて述べる。
 \item 原理と一致した点・一致しなかった点を示す。原因は確認済みの事実と推測を分ける。
 \item 実施範囲で言えることと、未確認の条件を一文ずつまとめる。
 \end{enumerate}
 \note{書き出しの型：「○○の条件で△△を測定したところ、□□となった。
 これは◇◇と整合する。ただし××は確認しておらず、一般化には追加測定が必要である。」
 空欄は自分の実測値と根拠で埋める。}
 \noindent 結論（実験後に記入）：\writing{30mm}
 \noindent 改善案・次に確認したい条件：\writing{16mm}
 \section*{参考文献}\label{sec:references}
 \begin{enumerate}
 \item \LabSource
 \item 使用したデータシート・書籍・ウェブ資料（著者、題名、該当箇所、URL、参照日）を追加する。
 \end{enumerate}
}
\newcommand{\prelabclosing}{
 \section*{準備・出典}\label{sec:prelabsource}
 \begin{itemize}[nosep]
 \item 資料、筆記具、記録用紙、電卓と、実験条件・機器設定の記録欄。
 \item 確認問題の解答と、分からない用語・回路点への具体的な質問。
 \end{itemize}
 {\small\noindent\LabSource\par}
}
\newcommand{\simpleflow}[1]{
 \begin{center}
 \begin{tikzpicture}[node distance=5mm and 5mm,
 every node/.style={font=\small},box/.style={draw=labblue,fill=labblue!4,rounded corners=1mm,
 align=center,minimum height=10mm,text width=30mm}]
 #1
 \end{tikzpicture}
 \end{center}
}

\newcommand{\LabCode}{C}
\newcommand{\LabTitle}{電源回路}
\newcommand{\SourcePdf}{IE3_expC_A4.pdf}
\begin{document}
\prelabtitle
\documentmap{
\tocrow{sec:auto1}{1 用語を一つずつ理解する}{言葉の意味、単位、回路や測定での役割を確認する。}
\tocrow{sec:auto2}{2 四つの回路の動きをたどる}{部品と信号の経路を追い、状態が変化する理由を理解する。}
\tocrow{sec:auto3}{3 整流と平滑の前提}{この項目の仕組みと実験で確認すべき点を整理する。}
\tocrow{sec:auto4}{4 小リプル近似の例題}{負荷・容量・充電間隔から脈動を見積もる。}
\tocrow{sec:auto5}{5 回路を読む練習}{部品と信号の経路を追い、状態が変化する理由を理解する。}
\tocrow{sec:auto6}{6 オシロスコープの読み方}{結合、倍率、ピーク間値、電流換算を確認する。}
\tocrow{sec:auto7}{7 測定からレポートへ：波形を数値にする}{実測値から計算・作図・記述へ進む順序を例題で練習する。}
\tocrow{sec:auto8}{8 確認問題}{自分で計算・説明して、実験に必要な理解を確かめる。}
\tocrow{sec:auto9}{解答の目安}{数値と理由を照合し、単位や論理の取り違えを見直す。}
\tocrow{sec:auto10}{実験前チェック}{接続・記録欄・必要な質問を準備する。}
\tocrow{sec:prelabsource}{準備・出典}{持参物と資料を確認し、具体的な質問を準備する。}
}

\section{用語を一つずつ理解する}\label{sec:auto1}
\begin{terms}
\term{整流}{交流から一定方向の出力を作る操作。整流直後の電圧はまだ大きく脈動しており、完全に一定の直流になったわけではない。}
\term{半波・全波}{半波は一方の半周期を使う。全波は両半周期を同じ向きへ変換するため、出力の山が一周期に二回現れる。}
\term{平滑}{整流後の脈動をCやLで小さくする操作。Cは電圧変化を、Lは電流変化を抑える性質を使う。}
\term{負荷電流}{電源から負荷へ供給する電流。大きいほど平滑コンデンサの放電が進みやすく、配線や素子での電圧降下も増える。}
\term{直流出力とリプル}{直流出力は平均的な電圧、リプルはその周囲の周期的変動。本資料のリプル量は最大と最小の差の半分である。}
\term{充電電流}{整流電圧がコンデンサ電圧を上回る区間で流れる電流。短い時間に電荷を補うため、瞬時値が平均負荷電流より大きくなることがある。}
\term{シャント抵抗}{電流測定用の小抵抗。両端電圧を測って$I=V/R$と換算する。小抵抗でも回路へ影響するので指定の値と接続点を使う。}
\end{terms}
\section{四つの回路の動きをたどる}\label{sec:auto2}
\subsection{C入力形：山で充電し、山の間で放電する}
ダイオードが導通できる区間にコンデンサへ電荷が補給される。
その後に整流電圧が下がるとダイオードは流れにくくなり、コンデンサから負荷へ電流が流れる。
負荷が重いほど一回の放電で減る電荷が大きい。出力波形の谷は低くなり、リプルが増えやすい。
全波では次の充電までの時間が半波より短いので、同じ条件ならリプルを小さくしやすい。
\subsection{チョーク入力形：電流を急に変えない}
整流直後の直列Lが電流の急変を抑え、その先のCが電圧の変動を抑える。
C入力形のように山付近の短い充電だけで供給する場合と比べ、電流の流れ方が異なる。
ただし軽負荷で電流が途切れると動作条件が変わるため、「どんな負荷でも同じ出力」とは言えない。
\clearpage
\subsection{全波ブリッジ：負荷の向きをそろえる}
入力が正の半周期では一組、負の半周期では別の組のダイオードが導通する。
どちらの経路でも負荷を流れる向きは同じである。負荷へ至る経路には通常二つのダイオードの順方向電圧降下が含まれる。
資料の回路図に、両半周期の導通経路を色分けして書き込むと理解しやすい。
\note{四条件の比較では負荷電流と電源条件をそろえる。回路の種類だけを変更した比較か、負荷も同時に変えた比較かを区別する。}

\section{整流と平滑の前提}\label{sec:auto3}
ダイオードは順方向に電流を流し、逆方向を遮断する。実素子には順方向電圧降下がある。
半波整流は交流の一方の半周期を使い、ブリッジ全波整流は両半周期を同じ極性の出力にする。
\[
 f_{r,\mathrm{half}}=f,\qquad f_{r,\mathrm{full}}=2f.
\]
C入力ではコンデンサが波形の山付近で充電され、山の間は負荷に放電する。
L入力ではインダクタが急な電流変化を抑える。
\[
 q=CV,\quad i_C=C\frac{dV_C}{dt},\quad v_L=L\frac{di_L}{dt}.
\]
\section{小リプル近似の例題}\label{sec:auto4}
放電中の電流を一定$I_L$とし、放電時間を$1/f_r$とすると
$\Delta V_{\mathrm{pp}}\simeq I_L/(f_rC)$となる。
$C=470\,\mu\mathrm{F},I_L=20\,\mathrm{mA},f=60\,\mathrm{Hz}$の例では、
\[
 \Delta V_{\mathrm{pp,half}}\simeq0.709\,\mathrm{V},\qquad
 \Delta V_{\mathrm{pp,full}}\simeq0.355\,\mathrm{V}.
\]
資料のリプル定義ではそれぞれ約$0.355\,\mathrm{V}$と$0.177\,\mathrm{V}$である。
\note{これは理論例である。負荷が重く、充電時間や出力低下が無視できない場合は、この近似だけで実験結果を説明できない。}
\section{回路を読む練習}\label{sec:auto5}
資料の四回路に、整流素子、平滑素子、負荷、直流測定点、リプル測定点、電流検出抵抗を書き込む。
\writing{35mm}
\clearpage
\section{オシロスコープの読み方}\label{sec:auto6}
\begin{itemize}
\item DC結合は直流成分を含む全波形を表示する。AC結合は直流成分を取り除くが、低周波では波形を変形させることもある。
\item ピーク間値は$V_{\max}-V_{\min}$。実効値や半ピーク間値と区別する。
\item 電流検出抵抗が$0.1\,\Omega$なら、$50\,\mathrm{mV}$は$0.50\,\mathrm{A}$に対応する。
\item 波形には時間/div、電圧/div、プローブ倍率、結合方式、観測点、トリガ源を記録する。
\end{itemize}
充電電流は平均負荷電流と等しい波形ではない。短時間の大きな充電電流が、次の放電までの電荷を補う。
\section{測定からレポートへ：波形を数値にする}\label{sec:auto7}
\subsection{リプルの定義をそろえる}
波形の最大と最小の差が4 div、設定が20 mV/div、未補正のプローブ倍率が10倍なら、
\[
 \Delta V_{\rm pp}=4\times20\,\mathrm{mV}\times10=0.80\,\mathrm{V},
 \qquad V_{\rm ripple}=0.40\,\mathrm{V}.
\]
これは読み取り練習の例。スコープが既に倍率を補正した表示なら、もう一度10倍してはいけない。
AC結合で直流成分が消えても、表示されたリプルの最大最小の差を読む方法は同じである。
\subsection{電流波形を読む}
$0.1\,\Omega$のシャントに、倍率補正後で0.10 Vが生じる瞬間は1.0 Aに対応する。
電流波形が短いパルスなら、高さだけで平均負荷電流と比べない。パルスの幅や繰返しも記録する。
$Q=\int i_Cdt$という面積の考え方で、放電して失った電荷をどの区間で補うかを説明できる。
\subsection{グラフと考察をつなぐ}
\begin{enumerate}
\item 横軸を負荷電流にそろえ、直流出力のグラフとリプルのグラフを別々に作る。
\item 各グラフに半波C、半波L、全波C、全波Lの四系列を同じ凡例で示す。
\item 一つの負荷で四条件を比べ、次に同じ回路で負荷増加の傾向を見る。
\item 500 mAと1 Aの波形を使い、充電の間隔・パルス幅・谷の深さの変化を説明する。
\end{enumerate}
小リプル近似は、放電中の電流と時間を単純化した見積もりである。
大きなリプルや軽負荷のチョーク回路へ、精密な予測式として無条件に使わない。

\clearpage
\clearpage
\section{確認問題}\label{sec:auto8}
\begin{enumerate}
\item 電源$60\,\mathrm{Hz}$で半波と全波のリプル周波数を答える。\answerline
\item 小リプル近似で負荷電流を2倍、容量を2倍にした場合をそれぞれ説明する。\answerline
\item $0.1\,\Omega$の検出抵抗に$0.12\,\mathrm{V}$が生じた瞬間の電流を求める。\answerline
\item 波形の最大値$12.8\,\mathrm{V}$、最小値$12.0\,\mathrm{V}$なら資料のリプル電圧はいくらか。\answerline
\end{enumerate}
\section*{解答の目安}\label{sec:auto9}
(1) $60\,\mathrm{Hz},120\,\mathrm{Hz}$。
(2) リプルはそれぞれ2倍、半分。
(3) $1.2\,\mathrm{A}$。
(4) $0.4\,\mathrm{V}$。
\section*{実験前チェック}\label{sec:auto10}
電源を切った状態で回路と接地を確認する。コンデンサの極性、耐圧、残留電荷に注意し、授業の装置操作手順に従う。
\prelabclosing
\end{document}
